\begin{exo}{}
	On considère la fonction $f$ définie sur $\intfo{-1}{+\infty}$ par :

	$$f(x) = \frac{6x + 6}{x + 2}$$

	\begin{enumerate}
		\item Pourquoi ne pouvait-on pas définir la fonction $f$ sur $\intfo{-2}{+\infty}$ ?
		\item Recopier et compléter le tableau de valeurs suivant.
	\end{enumerate}

	\begin{center}
		\def\arraystretch{1.4}
		\begin{tabular}{|c|*6{C{0.4cm}|}}
			\hline
			$x$    & $-1$ & $0$ & $1$ & $2$ & $4$ & $8$ \\ \hline
			$f(x)$ &      &     &     &     &     &     \\ \hline
		\end{tabular}
	\end{center}

	\begin{enumerate}[resume]
		\item Tracer la courbe représentative de $f$ dans un repère orthonormé.
		\item Dresser son tableau de variations sur $\intfo{-1}{+\infty}$ et préciser ses éventuels extremums.
	\end{enumerate}
\end{exo}
